% !TeX root = ../papers/nc-p3.tex
\section{有限点グラフの中心での倍率安定群と奇素数指数の排除}
\label{sec:odd-prime-finite-indices}
\begin{theorem}\label{thm:odd-prime-finite-indices}
中心有限な四辺有限の標準正則窓の一般頂点隅の可除指数は奇素数にならない。
より一般に、その指数 $s$ の奇素因子 $p$ に対しては、経路長4を次数1として作用する
倍率群の $\mu_p$ が固定する非原点の閉中心指標が存在する。
その指標の被約trace表現の全因子を、一つの薄い単純の $p$ 個の倍率だけで尽くすことはできない。
\end{theorem}
\begin{proof}
折り畳みの中心を $Z$、零番隅を $R$ とする。既出の中心有限な正則窓の構造により
$Z$ は次数零部分が $k$ の有限生成正規整域、$R$ のその分数体上の次元は $s^2$。
中心の次数 $n$ は経路長 $4n$ に対応する。一つの位置の矢だけを倍率変更すれば
この次数の通常の $\mathbb G_m$ 作用を得る。

$\mu_p$ が非原点を固定しなければ、各非原点中心指標では $p$ と素な次数の斉次元の
どれかが非零である。有限生成系のうちの有限個でこの条件を満たせる。
その次数の最小公倍数 $m$ は $p$ と素で、次数 $m$ に冪を揃えると原点外に共通零点がない。
従って一般の四元が次数 $m$ の多項式パラメータ環を与え、$Z,R$ はその上有限である。
$\dim R_n=\binom{4n+3}3$ からパラメータ環上の階数は $64m^4$ であり、
中心分数体を経由すれば $s^2$ の倍数である。奇素数 $p\mid s$ と矛盾する。

固定指標 $x$ で一般単純の被約traceを特殊化する。これは各頂点で総重み $s$ の
半単純型を与え、その繊維の全単純が正の重みで現れる
（定理\ref{thm:pi-vertex-extension-criterion}のtrace繊維の議論）。
非原点なので頂点単純はなく、各単純の四つの頂点次元は等しい。
中心指標が頂点と異なり全Extが消えるため、頂点とのEuler式の四階差が零だからである。
被約traceは代数自己同型と可換であり、$\mu_p$ はこの半単純型を重複度とともに保つ。

安定な単純を $S$ とする。一周に一回だけ矢を原始 $p$ 乗根だけ変える同型を取る。
その $p$ 乗はスカラーなので規格化して恒等とできる。固有空間に分ければ
$4p$ 個の位置を巡る表示となり、元の各頂点は $p$ 個の固有空間の和になる。
一箇所でも空なら全ての十分長い経路が零となり、元の非頂点単純性に反する。
よって頂点次元は少なくとも $p$ で、等号なら全固有空間が一次元になる。
間の矢も同時には零にならないので、これは長さ $4p$ の実際の周期点列である。
有限窓の標準四移動は点グラフ上で恒等だから、四つ後の各矢の射影座標は同じである。
従って元の表示で各位置の四行列は、一つの可逆行列のスカラー倍に限られる。
四位置のそれらの積の固有直線を順に運べば、頂点次元1の真部分表示ができ矛盾する。
安定な非頂点単純の頂点次元は厳密に $p$ より大きい。

$s=p$ なら従って安定因子はなく、重みを数えると薄い単純一つの $p$ 個の相異なる倍率だけが現れる。
同じグラフ点で一周倍率が異なる薄い単純の間の第一Extは両向きとも零である
（定理\ref{thm:one-corner-clifford-classification}の非対角矢の計算）。
これは同じGröbner表示や被約点の仮定を必要とせず、四移動の接写像が恒等であることによる。
中心有限代数の一つの完備局所中心上の閉繊維では、単純のExtグラフは連結である。
実際繊維代数で二群間に矢がなければ根基のフィルトレーションでブロックが分かれ、
完備環へ冪等元を持ち上げると両群間の隅は中山の補題で零となる。
中心は平坦な完備化と可換で局所環なのでこの分解はない。
これで薄い一軌道だけの繊維を排除し、$s=p$ の矛盾を得た。
\end{proof}
